\scanpage{40}
\mspage{39}
% Three oblique strokes cancel the first definition and stop above the fresh Soit.
\deleted{\uline{Définition} \quad Soit $(\deleted{\unsure{I|\mathcal J}}\unsure{\mathcal F|\mathscr F},\leq)$ un ens ordonné \inserted{avec $\mathrm{Sup}$ majorés et plus petit élément $\varnothing_{\mathcal F}$}. Une partie fermée $\mathcal F'$ de $\mathcal F$ est dite \uline{ample} si $\mathcal F'$ \deleted{\unsure{est|$\ne\varnothing$}} est stable par $\mathrm{Sup}$. majorés \uline{dans} $\mathcal F'$, \inserted{et $\mathcal F'\ne\varnothing$ $(\Longleftrightarrow\varnothing_{\mathcal F}\in\mathcal F')$.} On dit qu'elle \unsure{est|\deleted{est}} est \uline{\unsure{spacieuse|spacieux}} si elle est \deleted{stable par} $=$ \deleted{\unsure{spa|sp}}stable par $\mathrm{Sup}$ majorés dans $\mathcal F$.}

Soit $(\mathcal F,\leq)$ un ens. ordonné \deleted{stable par} \inserted{avec} $\mathrm{Sup}$. majorés et plus petit \unsure{élément|élmt} $\varnothing_{\mathcal F}$\inserted{(i.e. \unsure{satisfais.|satisf.} à At1}.

Soit $\deleted{\unsure{F|\mathcal J}}\mathcal F'$ une partie \uline{fermée}. Alors donc $\mathcal F'$ \deleted{\unsure{lors|les}} est non vide sss $\varnothing_{\mathcal F}\in\mathcal F'$. D'autre part, si $(F_i)$ est une famille dans $\mathcal F'$ qui est majorée dans $\mathcal F'$ par $\deleted{\unsure{F|H}}G$, alors son sup dans $\mathcal F$, soit $F$, est $\leq G$, donc $\in\mathcal F'$, et c'est un sup dans $\mathcal F'$. Donc $\mathcal F'$ est stable par sup. majorés dans $\mathcal F'$\inserted{(il satisfait donc At 1)}, et $\mathcal F'\hookrightarrow\mathcal F$ commute à ces sup.

Soit $\deleted{\unsure{X|N}}M$ l'ens des éléments (\unsure{str$^{\mathrm{t}}$|str.}) irréductibles de $\mathcal F$, et
\[
M'=M\cap\mathcal F'.
\]
Alors : $M'$ est l'ens des objets (str.) irréd. de $\mathcal F'$. En d'autres termes, si $F\in\mathcal F'$, pour que $F$ soit irréd. dans $\mathcal F'$, il f et s qu'il le soit dans $\mathcal F$. C'est immédiat, en termes des définitions (les familles $(\deleted{\unsure{d|A}}F_i)$ d'objets \deleted{\unsure{de|du}} majorés par $F$, que ce soit ds $\mathcal F$ ou $\mathcal F'$, sont les $=\mathrm{s}$, et leurs $\mathrm{Sup}$ dans $\mathcal F$ ou $\mathcal F'$ \unsure{itou|idem}, donc celles pour lesquelles $F=\mathrm{Sup}\,F_i$ sont les $=\mathrm{s}$, et dire que pour
